%% Abstract -- kept free of heavy mathematics per the IEEE template rule; all
%% headline numbers are emitted by the analysis pipeline as macros (never
%% hand-typed).
\begin{abstract}
Agentic AI increasingly writes the parallel code at the heart of high-performance
computing, and increasingly \emph{accepts its own work} through execution
feedback---the same deterministic verifier signal that verification-guided
reinforcement learning uses as its reward in place of human preference. In this
setting a program is accepted, and a post-training reward is paid, only when a
verifier confirms success, so the design of that verifier is the operational
definition of ``good'' parallel code. We ask what it must check. The natural gate is
configuration-robust correctness---agreement with a trusted serial reference under
\emph{every} thread count and repeated run---which catches the silent races a single
execution hides. We build an agent that writes OpenMP code and can invoke a
distribution-aware differential verifier establishing exactly this property, and
evaluate it on the ParEval benchmark with \Nmodels{} state-of-the-art models. Two
findings show a single-axis correctness gate is the wrong reward. First, it is
\emph{saturated}: single-shot generation is already robustly correct \Lzerorate{} of
the time, a repair loop raises this to \Lonerate{}, and replacing the loop's
single-thread check with the full differential verifier leaves it unchanged at
\Ltworate{}---these models seldom write the silent-race failure mode the verifier
targets, a null an independent happens-before detector (Archer) corroborates on every
accepted program. Second, and decisively, it is \emph{gameable by construction}: a
serial program with no parallelism agrees with the reference at every thread count and
scores a perfect \Ltworate{}-equivalent, so a reward defined on outputs is maximized
by deleting the parallelism HPC exists for. Measuring the second axis directly, the
\Ntimed{} programs the verifier accepted as robust span $7.9\times$ speedup down to
$0.33\times$---one accepted program runs three times \emph{slower} on eight threads
than one, and the correctness gate cannot tell it apart from a program that scales. We
argue that verification-guided post-training for HPC therefore needs a second,
\emph{performance} verification gate---does the accepted program actually scale---so
that the reward is two-axis, correctness \emph{and} realized parallelism; the first
has become cheap, and the second is where accepted code now fails. We release the
harness, the two-axis scorer, prompts, and complete agent transcripts.
\end{abstract}
